\documentclass[10pt,a4paper,landscape]{article}
\usepackage[utf8]{inputenc}
\usepackage[spanish,es-noshorthands]{babel}
\usepackage{tikz}
\usetikzlibrary{trees, positioning}
\usepackage[margin=1cm]{geometry}
\usepackage{amsmath}

\begin{document}

\begin{center}
    {\Large \textbf{Árbol Genealógico - Familia Rojas}} \\[10pt]
    \textit{Nota: Transcripción parcial. El árbol original fluye de derecha a izquierda.}
\end{center}

\vspace{1cm}

\begin{center}
\noindent\resizebox{\linewidth}{!}{%
\begin{tikzpicture}[
    grow=left,
    level distance=4.5cm,
    level 1/.style={sibling distance=8cm},
    level 2/.style={sibling distance=3cm},
    level 3/.style={sibling distance=1.2cm},
    edge from parent path={(\tikzparentnode.west) -- +( -0.5cm, 0) |- (\tikzchildnode.east)},
    every node/.style={draw, rectangle, rounded corners, font=\small, align=center, fill=white},
    anchor=east
]

\node {Baltazar \\ Rojas}
    child { node {María del Rosario \\ Rojas}
        child { node {Andrés \\ Marín Rojas} 
            child { node {Flor María} }
            child { node {Deyanira} }
            child { node {Juan \\ Antonio} }
            child { node {Miguel \\ Antonio} }
            child { node {José Miguel} }
            child { node {Cruz} }
        }
        child { node {Alejandrina \\ Marín T.} 
            child { node {Josefina} }
            child { node {Argentina} }
            child { node {Ema} }
            child { node {Rosa \\ Pastora} }
            child { node {Sara Rosa} }
        }
        child { node {Domingo \\ Marín T.} 
            child { node {Marta} }
            child { node {Rosario} }
            child { node {María Elena} }
            child { node {Elisa} }
        }
        child { node {Margarita \\ Marín Rojas} 
            child { node {Rafael} }
        }
    };
\end{tikzpicture}%
}
\end{center}

\end{document}
